\documentclass{amsart}
\usepackage[UTF8,fontset=windows]{ctex}
\usepackage[a4paper,margin=24mm]{geometry}
% For dealing with references we use the comment environment
\usepackage{verbatim}
\newenvironment{reference}{\comment}{\endcomment}
%\newenvironment{reference}{}{}
\newenvironment{slogan}{\comment}{\endcomment}
\newenvironment{history}{\comment}{\endcomment}

% For commutative diagrams we use Xy-pic
\usepackage[all]{xy}

% We use 2cell for 2-commutative diagrams.
\xyoption{2cell}
\UseAllTwocells

% We use multicol for the list of chapters between chapters
\usepackage{multicol}

% This is generally recommended for better output
\usepackage{lmodern}
\usepackage[T1]{fontenc}

% For cross-file-references

% Package for hypertext links:
\usepackage{hyperref}

% For any local file, say "hello.tex" you want to link to please

% Theorem environments.
%
\theoremstyle{plain}
\newtheorem{theorem}[subsection]{Theorem}
\newtheorem{proposition}[subsection]{Proposition}
\newtheorem{lemma}[subsection]{Lemma}

\theoremstyle{definition}
\newtheorem{definition}[subsection]{Definition}
\newtheorem{example}[subsection]{Example}
\newtheorem{exercise}[subsection]{Exercise}
\newtheorem{situation}[subsection]{Situation}

\theoremstyle{remark}
\newtheorem{remark}[subsection]{Remark}
\newtheorem{remarks}[subsection]{Remarks}

\numberwithin{equation}{subsection}

% Macros
%
\def\lim{\mathop{\mathrm{lim}}\nolimits}
\def\colim{\mathop{\mathrm{colim}}\nolimits}
\def\Spec{\mathop{\mathrm{Spec}}}
\def\Hom{\mathop{\mathrm{Hom}}\nolimits}
\def\Ext{\mathop{\mathrm{Ext}}\nolimits}
\def\SheafHom{\mathop{\mathcal{H}\!\mathit{om}}\nolimits}
\def\SheafExt{\mathop{\mathcal{E}\!\mathit{xt}}\nolimits}
\def\Sch{\mathit{Sch}}
\def\Mor{\mathop{\mathrm{Mor}}\nolimits}
\def\Ob{\mathop{\mathrm{Ob}}\nolimits}
\def\Sh{\mathop{\mathit{Sh}}\nolimits}
\def\NL{\mathop{N\!L}\nolimits}
\def\CH{\mathop{\mathrm{CH}}\nolimits}
\def\proetale{{pro\text{-}\acute{e}tale}}
\def\etale{{\acute{e}tale}}
\def\QCoh{\mathit{QCoh}}
\def\Ker{\mathop{\mathrm{Ker}}}
\def\Im{\mathop{\mathrm{Im}}}
\def\Coker{\mathop{\mathrm{Coker}}}
\def\Coim{\mathop{\mathrm{Coim}}}

% Boxtimes
%
\DeclareMathSymbol{\boxtimes}{\mathbin}{AMSa}{"02}

%
% Macros for moduli stacks/spaces
%
\def\QCohstack{\mathcal{QC}\!\mathit{oh}}
\def\Cohstack{\mathcal{C}\!\mathit{oh}}
\def\Spacesstack{\mathcal{S}\!\mathit{paces}}
\def\Quotfunctor{\mathrm{Quot}}
\def\Hilbfunctor{\mathrm{Hilb}}
\def\Curvesstack{\mathcal{C}\!\mathit{urves}}
\def\Polarizedstack{\mathcal{P}\!\mathit{olarized}}
\def\Complexesstack{\mathcal{C}\!\mathit{omplexes}}
% \Pic is the operator that assigns to X its picard group, usage \Pic(X)
% \Picardstack_{X/B} denotes the Picard stack of X over B
% \Picardfunctor_{X/B} denotes the Picard functor of X over B
\def\Pic{\mathop{\mathrm{Pic}}\nolimits}
\def\Picardstack{\mathcal{P}\!\mathit{ic}}
\def\Picardfunctor{\mathrm{Pic}}
\def\Deformationcategory{\mathcal{D}\!\mathit{ef}}

\hypersetup{hidelinks}
\widowpenalty=10000\clubpenalty=10000\displaywidowpenalty=10000
\begin{document}
\title[Stacks Project, Tag 0BKI]{任意滤过复形的 K-injective 兼容解消}
\maketitle
\noindent\textbf{Editorial reading note.} Author text is retained verbatim. In the product over p-prime greater than p, the factor printed J\textasciicircum{}\{p,bullet\} is indexed by p-prime as in the surrounding displayed definitions. The composition on source line 2214 has the order of the two alpha factors reversed by type. Quotients F\textasciicircum{}p/F\textasciicircum{}\{p-prime\} use p-prime at least p. These are editorial reading notes, not changes to the author body.


\noindent\textbf{Counting scope.} The primary problem is Tag 0BKI, marked below. All preceding numbered statements are supporting author context and are not counted as additional problems. This is a local pilot with a temporary ID, not production admission.


\noindent Copyright (C) 2005--2025 Johan de Jong. Stacks Project authors. Source: \href{https://stacks.math.columbia.edu/tag/0BKI}{Tag 0BKI}.

\noindent Difficulty H2: 本题要求同时控制复形、各滤过层及所有相关商的逐项内射性和K-injective性。单独解消各层不会给出单射相容滤过；作者用无环K-injective辅助复形及无限乘积构造新的滤过，再以逐项分裂正合列同时验证所有商与拟同构。允许未滤过解消定理后仍有新的相容构造，判H2。.

\noindent History: excerpt from revision local cached snapshot; exact chapter hashes; revision unknown, injectives.tex, lines 2107--2240. Only reference presentation was adapted. Licensed under GNU FDL 1.2 or later; no invariant sections or cover texts. See accompanying COPYING.

\begin{definition}
\label{homology-definition-filtered-complex}
Let $\mathcal{A}$ be an abelian category.
A {\it filtered complex $K^\bullet$ of $\mathcal{A}$}
is a complex of $\text{Fil}(\mathcal{A})$ (see
Definition \ref{homology-definition-filtered}).
\end{definition}

\noindent
We will denote the filtration on the objects by $F$. Thus
$F^pK^n$ denotes the $p$th step in the filtration of the $n$th term of the
complex. Note that each $F^pK^\bullet$ is a complex of $\mathcal{A}$.
Hence we could also have defined a filtered complex as a filtered object
in the (abelian) category of complexes of $\mathcal{A}$.
In particular $\text{gr} K^\bullet$ is a graded object of the
category of complexes of $\mathcal{A}$.


\begin{definition}
\label{homology-definition-filtered}
Let $\mathcal{A}$ be an abelian category.
\begin{enumerate}
\item A {\it decreasing filtration} $F$ on an object $A$
is a family $(F^nA)_{n \in \mathbf{Z}}$ of subobjects of $A$ such that
$$
A \supset \ldots \supset F^nA \supset F^{n + 1}A \supset \ldots \supset 0
$$
\item A {\it filtered object of $\mathcal{A}$} is
pair $(A, F)$ consisting of an object $A$ of $\mathcal{A}$
and a decreasing filtration $F$ on $A$.
\item A {\it morphism $(A, F) \to (B, F)$ of filtered objects}
is given by a morphism $\varphi : A \to B$ of $\mathcal{A}$
such that $\varphi(F^iA) \subset F^iB$ for all $i \in \mathbf{Z}$.
\item The category of filtered objects is denoted $\text{Fil}(\mathcal{A})$.
\item Given a filtered object $(A, F)$ and a subobject $X \subset A$ the
{\it induced filtration} on $X$ is the filtration with $F^nX = X \cap F^nA$.
\item Given a filtered object $(A, F)$ and a surjection
$\pi : A \to Y$ the {\it quotient filtration} is the filtration with
$F^nY = \pi(F^nA)$.
\item A filtration $F$ on an object $A$ is said to be {\it finite}
if there exist $n, m$ such that $F^nA = A$ and $F^mA = 0$.
\item Given a filtered object $(A, F)$ we say $\bigcap F^iA$ exists
if there exists a biggest subobject of $A$ contained in all $F^iA$.
We say $\bigcup F^iA$ exists if there exists a smallest subobject
of $A$ containing all $F^iA$.
\item The filtration on a filtered object $(A, F)$ is said to be
{\it separated} if $\bigcap F^iA = 0$ and
{\it exhaustive} if $\bigcup F^iA = A$.
\end{enumerate}
\end{definition}

\begin{definition}
\label{derived-definition-K-injective}
Let $\mathcal{A}$ be an abelian category. A complex $I^\bullet$
is {\it K-injective} if for every acyclic complex $M^\bullet$ we
have $\Hom_{K(\mathcal{A})}(M^\bullet, I^\bullet) = 0$.
\end{definition}

\begin{theorem}
\label{theorem-K-injective-embedding-grothendieck}
\begin{slogan}
Existence of K-injective complexes for Grothendieck abelian categories.
\end{slogan}
Let $\mathcal{A}$ be a Grothendieck abelian category.
For every complex $M^\bullet$ there exists a quasi-isomorphism
$M^\bullet \to I^\bullet$ such that $M^n \to I^n$ is injective and $I^n$
is an injective object of $\mathcal{A}$ for all $n$ and $I^\bullet$
is a K-injective complex. Moreover, the construction is functorial in
$M^\bullet$.
\end{theorem}

\begin{lemma}
\label{lemma-derived-products}
Let $\mathcal{A}$ be a Grothendieck abelian category.
Then
\begin{enumerate}
\item $D(\mathcal{A})$ has both direct sums and products,
\item direct sums are obtained by taking termwise direct sums of
any complexes,
\item products are obtained by taking termwise products of
K-injective complexes.
\end{enumerate}
\end{lemma}

\begin{proof}
Let $K^\bullet_i$, $i \in I$ be a family of objects of $D(\mathcal{A})$
indexed by a set $I$. We claim that the termwise direct sum
$\bigoplus_{i \in I} K^\bullet_i$ is a direct sum in $D(\mathcal{A})$.
Namely, let $I^\bullet$ be a K-injective complex. Then we have
\begin{align*}
\Hom_{D(\mathcal{A})}(\bigoplus\nolimits_{i \in I} K^\bullet_i, I^\bullet)
& =
\Hom_{K(\mathcal{A})}(\bigoplus\nolimits_{i \in I} K^\bullet_i, I^\bullet) \\
& =
\prod\nolimits_{i \in I} \Hom_{K(\mathcal{A})}(K^\bullet_i, I^\bullet) \\
& =
\prod\nolimits_{i \in I} \Hom_{D(\mathcal{A})}(K^\bullet_i, I^\bullet)
\end{align*}
as desired. This is sufficient since any complex can be represented
by a K-injective complex by
Theorem \ref{theorem-K-injective-embedding-grothendieck}.
To construct the product, choose a K-injective resolution
$K_i^\bullet \to I_i^\bullet$ for each $i$. Then we claim that
$\prod_{i \in I} I_i^\bullet$ is a product in $D(\mathcal{A})$.
This follows from
Derived Categories, Lemma \href{https://stacks.math.columbia.edu/tag/0BK6}{lemma product K injective (Tag 0BK6)}.
\end{proof}


\begin{lemma}
\label{derived-lemma-triangle-K-injective}
Let $\mathcal{A}$ be an abelian category. Let $(K, L, M, f, g, h)$
be a distinguished triangle of $K(\mathcal{A})$. If two out of
$K$, $L$, $M$ are K-injective complexes, then the third is too.
\end{lemma}

\section*{Primary problem: Tag 0BKI}
\begin{lemma}
\label{lemma-K-injective-embedding-filtration}
Let $\mathcal{A}$ be a Grothendieck abelian category.
Let $K^\bullet$ be a filtered complex of $\mathcal{A}$, see
Homology, Definition \ref{homology-definition-filtered-complex}.
Then there exists a morphism $j : K^\bullet \to J^\bullet$
of filtered complexes of $\mathcal{A}$ such that
\begin{enumerate}
\item $J^n$, $F^pJ^n$, $J^n/F^pJ^n$ and $F^pJ^n/F^{p'}J^n$ are injective
objects of $\mathcal{A}$,
\item $J^\bullet$, $F^pJ^\bullet$, $J^\bullet/F^pJ^\bullet$, and
$F^pJ^\bullet/F^{p'}J^\bullet$ are K-injective complexes,
\item $j$ induces quasi-isomorphisms
$K^\bullet \to J^\bullet$,
$F^pK^\bullet \to F^pJ^\bullet$,
$K^\bullet/F^pK^\bullet \to J^\bullet/F^pJ^\bullet$, and
$F^pK^\bullet/F^{p'}K^\bullet \to F^pJ^\bullet/F^{p'}J^\bullet$.
\end{enumerate}
\end{lemma}

\begin{proof}
By Theorem \ref{theorem-K-injective-embedding-grothendieck}
we obtain quasi-isomorphisms $i : K^\bullet \to I^\bullet$ and
$i^p : F^pK^\bullet \to I^{p, \bullet}$ as well as commutative diagrams
$$
\vcenter{
\xymatrix{
K^\bullet \ar[d]_i & F^pK^\bullet \ar[l] \ar[d]_{i^p} \\
I^\bullet & I^{p, \bullet} \ar[l]_{\alpha^p}
}
}
\quad\text{and}\quad
\vcenter{
\xymatrix{
F^{p'}K^\bullet \ar[d]_{i^{p'}} &
F^pK^\bullet \ar[l] \ar[d]_{i^p} \\
I^{p', \bullet} &
I^{p, \bullet} \ar[l]_{\alpha^{p p'}}
}
}
\quad\text{for }p' \leq p
$$
such that $\alpha^p \circ \alpha^{p' p} = \alpha^{p'}$
and $\alpha^{p'p''} \circ \alpha^{pp'} = \alpha^{pp''}$.
The problem is that the maps $\alpha^p : I^{p, \bullet} \to I^\bullet$
need not be injective. For each $p$ we choose an injection
$t^p : I^{p, \bullet} \to J^{p, \bullet}$ into an acyclic K-injective
complex $J^{p, \bullet}$ whose terms are injective objects of $\mathcal{A}$
(first map to the cone on the identity and then use the theorem).
Choose a map of complexes $s^p : I^\bullet \to J^{p, \bullet}$
such that the following diagram commutes
$$
\xymatrix{
K^\bullet \ar[d]_i & F^pK^\bullet \ar[l] \ar[d]_{i^p} \\
I^\bullet \ar[rd]_{s^p} & I^{p, \bullet} \ar[d]^{t^p} \\
& J^{p, \bullet}
}
$$
This is possible: the composition $F^pK^\bullet \to J^{p, \bullet}$
is homotopic to zero because $J^{p, \bullet}$ is acyclic and K-injective
(Derived Categories, Lemma \href{https://stacks.math.columbia.edu/tag/070I}{lemma K injective (Tag 070I)}).
Since the objects $J^{p, n - 1}$ are injective and since
$F^pK^n \to K^n \to I^n$ are injective morphisms, we
can lift the maps $F^pK^n \to J^{p, n - 1}$ giving the homotopy
to a map $h^n : I^n \to J^{p, n - 1}$. Then we set $s^p$
equal to $h \circ \text{d} + \text{d} \circ h$.
(Warning: It will not be the case that $t^p = s^p \circ \alpha^p$,
so we have to be careful not to use this below.)

\medskip\noindent
Consider
$$
J^\bullet = I^\bullet \times \prod\nolimits_p J^{p, \bullet}
$$
Because products in $D(\mathcal{A})$ are given by taking
products of K-injective complexes
(Lemma \ref{lemma-derived-products})
and since $J^{p, \bullet}$
is isomorphic to $0$ in $D(\mathcal{A})$ we see that
$J^\bullet \to I^\bullet$ is an isomorphism in $D(\mathcal{A})$.
Consider the map
$$
j = i \times (s^p \circ i)_{p \in \mathbf{Z}} :
K^\bullet
\longrightarrow
I^\bullet \times \prod\nolimits_p J^{p, \bullet} = J^\bullet
$$
By our remarks above this is a quasi-isomorphism. It is also injective.
For $p \in \mathbf{Z}$ we let $F^pJ^\bullet \subset J^\bullet$ be
$$
\Im\left(
\alpha^p \times (t^{p'} \circ \alpha^{pp'})_{p' \leq p} :
I^{p, \bullet}
\to
I^\bullet \times \prod\nolimits_{p' \leq p} J^{p', \bullet}
\right)
\times \prod\nolimits_{p' > p} J^{p', \bullet}
$$
This complex is isomorphic to the complex
$I^{p, \bullet} \times \prod_{p' > p} J^{p, \bullet}$
as $\alpha^{pp} = \text{id}$ and $t^p$ is injective.
Hence $F^pJ^\bullet$ is quasi-isomorphic to $I^{p, \bullet}$ (argue
as above). We have $j(F^pK^\bullet) \subset F^pJ^\bullet$ because
of the commutativity of the diagram above. The corresponding
map of complexes $F^pK^\bullet \to F^pJ^\bullet$ is a quasi-isomorphism
by what we just said. Finally, to see that
$F^{p + 1}J^\bullet \subset F^pJ^\bullet$
use that $\alpha^{p + 1p} \circ \alpha^{pp'} = \alpha^{p + 1p'}$
and the commutativity of the first displayed diagram
in the first paragraph of the proof.

\medskip\noindent
We claim that $j : K^\bullet \to J^\bullet$ is a solution to the
problem posed by the lemma. Namely, $F^pJ^n$ is an injective object
of $\mathcal{A}$ because it is isomorphic to
$I^{p, n} \times \prod_{p' > p} J^{p', n}$ and products of
injectives are injective. Then the injective map $F^pJ^n \to J^n$
splits and hence the quotient $J^n/F^pJ^n$ is injective as well
as a direct summand of the injective object $J^n$.
Similarly for $F^pJ^n/F^{p'}J^n$. This in particular means
that $0 \to F^pJ^\bullet \to J^\bullet \to J^\bullet/F^pJ^\bullet \to 0$
is a termwise split short exact sequence of complexes, hence defines
a distinguished triangle in $K(\mathcal{A})$ by fiat.
Since $J^\bullet$ and $F^pJ^\bullet$ are K-injective complexes
we see that the same is true for $J^\bullet/F^pJ^\bullet$
by Derived Categories, Lemma \ref{derived-lemma-triangle-K-injective}.
A similar argument shows that $F^pJ^\bullet/F^{p'}J^\bullet$
is K-injective. By construction $j : K^\bullet \to J^\bullet$
and the induced maps $F^pK^\bullet \to F^pJ^\bullet$ are
quasi-isomorphisms. Using the long exact cohomology sequences
of the complexes in play we find that the same holds for
$K^\bullet/F^pK^\bullet \to J^\bullet/F^pJ^\bullet$ and
$F^pK^\bullet/F^{p'}K^\bullet \to F^pJ^\bullet/F^{p'}J^\bullet$.
\end{proof}
\end{document}
